\section{Discussion: each axis bites at a different stage}\label{sec:discussion}

Table~\ref{tab:staged} re-estimates every outcome on one harmonized
sample --- one row per debut owner, scored on that owner's first filing
--- under a single specification, so the coefficients can be read down a
column. Magnitudes differ from the section-specific tables above.

\begin{table}[h]\centering\small
\caption{Every outcome in the funnel, conditioned on its prerequisite, under one specification. Unit is the debut owner, scored on that owner's first filing. Each row is a linear probability model with class$\times$debut-year fixed effects and heteroskedasticity-robust errors; coefficients are percentage points per standard deviation. \emph{Atypicality} is the average of the two KL levels; \emph{lead} is the signed difference. Cohort windows differ because each maintenance gate needs elapsed time and Form~D is electronic only from 2009, so the rows are not a nested decomposition. The first gate is a dated cancellation for non-use; the second is whether a registration that cleared the first was renewed rather than cancelled or expired, restricted to 2002--2013 cohorts, since a year-six and a year-ten death share a terminal status and are separable only this way.}
\label{tab:staged}
\begin{tabular}{llrrrrr}
\toprule
 & & & & \multicolumn{2}{c}{Unsigned} & Signed \\
\cmidrule(lr){5-6}\cmidrule(lr){7-7}
Outcome & Given & $n$ & Base & Atypicality & $|\Delta$KL$|$ & Lead \\
\midrule
Reached registration & filed & 3,161,403 & 56.12\% & $-0.903$ (0.032) & $-0.344$ (0.047) & $+1.095$ (0.029) \\
Passed first gate (yr 6) & registered & 1,774,111 & 47.94\% & $+0.858$ (0.043) & $-0.449$ (0.062) & $-0.955$ (0.038) \\
Passed second gate (yr 10) & passed first & 298,002 & 55.92\% & $+0.705$ (0.102) & $-0.294$ (0.171) & $-0.861$ (0.104) \\
Raised a Reg D round & filed & 1,534,824 & 2.53\% & $-0.023$ (0.015) & $+0.034$ (0.030) & $-0.272$ (0.018) \\
Raised a Reg D round & registered & 916,442 & 2.61\% & $-0.056$ (0.020) & $-0.015$ (0.040) & $-0.235$ (0.023) \\
Listed after the round & funded & 38,886 & 3.49\% & $+0.445$ (0.123) & $+0.618$ (0.208) & $-0.641$ (0.128) \\
Reached listed universe & registered & 1,774,111 & 0.23\% & $+0.030$ (0.005) & $+0.008$ (0.008) & $-0.002$ (0.005) \\
\bottomrule
\end{tabular}
\\[0.3em]\footnotesize Standard errors in parentheses, in place of significance stars; see the note to Table~\ref{tab:gate_lpm}.
\end{table}

Surprising language costs at registration and pays afterwards: atypical
debut filings complete registration less often ($-1.35$pp per standard
deviation), then survive the five-year proof of continued use more often
($+1.35$pp) and reach public markets more often ($+0.04$pp on a 0.56\%
base --- the one specification in which a positive atypicality
coefficient on listing survives, and it is a continuous term through a
non-monotone profile, not the quintile contrast of
Table~\ref{tab:edgar}). Having been early runs the other way: leading
applications complete registration \emph{more} often ($+0.87$pp) and then
pass the five-year proof \emph{less} often ($-0.49$pp on passing). Both
survival rows point the same way at the year-ten renewal, four years
later and on a sixth of the sample: atypicality $+1.30$pp and lead
$-0.63$pp among marks that had already cleared the five-year proof.
Whatever the five-year proof selects on, the renewal selects on it again
rather than reversing it. The registration coefficient on lead is linear
through a profile that is nearly flat (Section~\ref{sec:reg_rates}); the
positive term reflects the asymmetry of the atypicality curve rather
than a monotone gain from leading.

In the financing rows atypicality does nothing at one stage and a great
deal at the next. Whether a debut owner ever raises a Regulation~D round
--- a private securities placement, reported to the SEC on Form~D --- is
unrelated to how atypical its filing was: $-0.034$pp per standard
deviation among owners who completed registration, against a 2.45\% base
($t = -1.7$). Among owners who did raise one, the same measure predicts
\emph{reaching} the public market: $+0.61$pp on a 3.51\% base, about 17\%
relative ($t = 4.4$). Atypical language is invisible at the financing stage and
informative afterwards, which is what an investor screen that does not
read the text would look like --- the text being free and public at the
moment of filing. The comparison is among survivors, so it cannot
separate that reading from selection on unobservables; and the null at
entry is a null, not evidence of indifference. One construction matters:
an exchange-registration marker on its own does not mean a firm listed
\emph{after} raising --- 55\% of matched owners carrying one registered
before their first Regulation~D notice, because established public
companies also place securities privately --- so the listing-after-round
row is restricted to owners whose exchange registration postdates their
first notice, which halves the coefficient and leaves it significant.

Three limits bound the financing rows. Regulation~D notices are
electronically filed only from 2009, so they rest on 2009--2018 debut
cohorts. Owners are matched by normalized name, which succeeds for a
small and unrepresentative minority favouring formally constituted
entities, so every estimate is conditional on matchability. And the final
row conditions on funding, realized after the filing being scored, so the
coefficient is a contrast among survivors rather than a treatment effect:
it establishes that information distinguishing eventual listers was
already present in the text years earlier, not that the language caused
the outcome.

Lead does not reverse anywhere after registration. It is negative at both
use proofs, at financing with or without registration imposed, and at
listing among the funded; the only row where it is not negative is SEC
reporting among registered owners, where it is indistinguishable from
zero ($-0.004$pp, SE $0.007$). Leading filings fare worse at every stage
where the market, rather than the examiner, is doing the selecting, and
are never rewarded at any of them.

Why would having been early cost? The two axes are not independent in the
obvious economic sense. What makes a filing leading is precisely that its
industry moved toward its language, and an industry that has moved toward
your language now describes its goods the way you describe yours. Being
early therefore consumes the quantity that pays: the mark that was
unusual in 2011 is ordinary by 2016 not because its owner changed
anything but because the category arrived. On this reading atypicality is
a position competitors erode, and lead measures how fast. The design
cannot demonstrate that mechanism, since it cannot separate a category
catching up with a filing from some third thing moving both. What can be
said is weaker: a negative coefficient on lead at equal atypicality is
what the mechanism would produce if it were real.

A second reading follows from the first, and the design cannot test it
either. Privately, writing ahead of the industry's language is hard to
rationalize: no stage of the record rewards it, and each era's leading
vocabulary drew filers by the thousands anyway. Socially, the
exploration is not obviously waste. The corpus moves --- what makes a
filing leading is precisely that its class later wrote the way it did
--- so the language every industry eventually standardized on was first
filed by entrants who disproportionately did not survive to use it. If
finding the vocabulary an industry will adopt has value to the eventual
adopters, that value accrued to the followers while the cost fell on the
explorers.

Read together, the record's stages sort the two properties cleanly. The
examiner rewards middling conformity and cannot see lead at all. The
market's first proof, five years on, punishes having been early --- most
where the language was moving fastest --- and spares the unusual, though
mostly the unusual that extends an established theme into a new class.
The financing market reads none of it at entry. And the firms that came
to hold most of the linked economy's value filed first descriptions in
the words their industries already used. The pioneers of language are the
ones with arrows in their backs; the settlers arrive writing prose
everyone recognises, and they are few.
